\subsection{Methods details}\label{sec:n-methods}
\vsette{Assembly identifiers, contigs, bounds and target-specific decisions are retained in \nolinkurl{data/markers_u/truth_regions.tsv}, \nolinkurl{data/method/dyz19_clean.tsv}, \nolinkurl{data/method/dyz17_clean.tsv} and \nolinkurl{data/method/truth_qc/x/rule_w4.txt}; \nolinkurl{results/figures/reference_targets.tsv} lists the target definitions before the final exclusions. Published Yq12 annotations are archived in \nolinkurl{data/method/yq12_truth_hallast.tsv}. The downloaded annotation release was HGSVC3 chrY analysis 1.0, in \nolinkurl{20260513_chrY_annotations/verkko-v2.2.1}. Extended Data Fig.~\ref{fig:array-structure} shows the arrays of the 156 assemblies in these panels. Inclusion in a pangenome resource or a favourable genome-wide quality metric did not replace array-specific assessment; repetitive Yq remains an assembly limitation in HPRC release 2 \cite{Lucas2026}.}

Estimates showed no clear length-dependent bias: slopes of log estimate on log reference length, one value per male, were 1.000 for DYZ3 (bootstrap 95\% interval 0.985--1.013), 1.008 for Yq12, 0.997 for DYZ1, 1.025 for DYZ2 (0.999--1.051), 0.970 for DYZ19, 0.949 for DYZ17 and 0.899 for DYZ18S (0.734--1.056). Absolute relative error is $100|\widehat L_i/L_i-1|$, and signed relative error is $100(\widehat L_i/L_i-1)$. Figure~\ref{fig:method}b gives each individual equal weight: each point is the arithmetic mean of ten absolute out-of-fold percentage errors, one per repeat. Errors are made absolute before averaging. Boxes span the 25th--75th percentiles, central lines show medians and whiskers end at the 10th and 90th percentiles; all individual points remain visible. Labels identify the three highest individual mean errors per target, with ties resolved by sample identifier. In the coverage panel, each point likewise averages ten absolute errors for one individual at one depth fraction. Curves connect the medians of these individual means. Boxes at different coverages are equally spaced and labelled by measured median autosomal coverage, rather than by the thinning fraction. Error axes are linear through 5\% and logarithmic above. Medians of the 50 test folds ranged from 1.48\% to 3.37\% for DYZ3 and from 0.77\% to 2.45\% for Yq12. Table~\ref{tab:accuracy} reports the median of all out-of-fold errors, which differs from the median of the displayed individual means; maxima refer to single out-of-fold predictions. Repeated predictions are not independent biological individuals.

The historical Miga/Altemose samples \cite{Miga2014,Altemose2014} informed method development. \vsette{With an earlier model} their median errors were 2.85\% for DYZ3 (20 individuals), 1.16\% for Yq12 (12) and 2.72\% for DYZ19 (21). These comparisons are treated as development evidence. The accuracy assessment reported here uses cross-validation; a separate cohort that did not inform method design is needed to evaluate external performance. \vsette{The two annotations agree closely, which limits differences between the two procedures but not errors they share or errors of the assemblies: against our profile-HMM annotation of the same 91 assemblies, the published lengths differ by a constant offset ($+1.92\%$ for DYZ1, $-0.27\%$ for DYZ2, $+0.05\%$ for DYZ18) that the scale absorbs, and around it by medians of 0.09\%, 0.07\% and 0.00\%. With exact marker copy counts of the assemblies in place of reads, the CV error of reference over marker copies was 0.60\% for DYZ1, 0.82\% for DYZ2 and 3.97\% for DYZ18S, so read and library effects dominate for DYZ1 and DYZ2 and marker density for DYZ18S.} \vsette{Empirical 90\% half-widths of the CV errors, as a percentage of the estimate, were 5.17\% for DYZ3, 4.23\% for Yq12, 5.45\% for DYZ1, 6.40\% for DYZ2, 10.69\% for DYZ18S, 6.76\% for DYZ19 and 6.11\% for DYZ17; the held-out males fell within them in 17 of 18, 3 of 3, 3 of 3, 3 of 3, 3 of 3, 14 of 14 and 15 of 15 cases. These are pooled quantiles of repeated out-of-fold predictions, not calibrated intervals for one sample. Finite-sample 90\% quantiles of one score per male (median absolute log error over the ten repeats) were 5.23\%, 4.49\%, 5.83\%, 6.86\%, 11.81\%, 6.74\% and 6.22\% in the same order, with the same held-out counts. No separate calibration set was used, so these widths are empirical summaries without a coverage guarantee. Three males cannot establish coverage for Yq12, and pooled widths hide specific states: the errors of three of four A males exceeded the DYZ18S width in all ten repeats ($-17.63\%$ to $-13.25\%$; the DYZ18S arrays of two of them carry 85.75\% of the DYZ18S markers, against a median of 97.08\% in 91 assemblies, 88.3\% of that median, in line with the underestimate; HG02984, the A0 male without the bias, carries 71.08\%, so presence does not explain every case, and across 76 males presence and CV error correlate at $-0.144$), HG02492 exceeded the DYZ19 width in nine of ten ($+6.88\%$ to $+25.47\%$) and three of six J males exceeded the Yq12 width in all ten ($-6.42\%$ to $-4.14\%$). These widths apply to PCR-free libraries of the calibration project only.}

\subsubsection{Pedigrees, relatedness and sex labels}
All \vsette{seven} satellite targets and the DYZ1:DYZ2 ratio are compared in Extended Data Fig.~\ref{fig:all-pedigrees}. Its signed-difference panels additionally retain the two reverse-direction pairs as a separate group. Signed differences are $100(\widehat L_{\mathrm{offspring}}/\widehat L_{\mathrm{parent}}-1)$, with absolute values used for concordance. Absolute-length axes have matching limits and a diagonal denoting equality. Shapes distinguish batch contrasts, including the reverse direction in the absolute-length panels, and open points mark pairs involving a male with a reference assembly. The largest labelled differences give offspring identifiers. All pairs contribute to overall concordance; the two reverse-direction pairs are not pooled into the signed-difference summaries of the main cross-batch direction. Signed-difference axes are linear within $\pm5\%$ and logarithmic outside; summary lines show medians and interquartile ranges. No batch adjustment was applied. Pedigree agreement combines measurement differences and potential biological variation and is not a mutation-rate estimate.

Close relatedness was assessed using two public GRCh38 KING-robust first/second-degree pair reports for the 3,202-person 1kGP collection from \url{https://www.cog-genomics.org/plink/2.0/resources} (retrieved 27 October 2026). Published estimates were joined to sample and pedigree records by exact identifiers, retaining both reported coefficients and IBS0 values. Both reports support all 309 parent--offspring pairs. \vsette{All 115} unique CV individuals are covered by the corrected public metadata; neither report adds a cross-FamilyID first/second-degree link among them. Ten pairs across different FamilyIDs occur in both reports, including one with kinship 0.264; two further pairs occur only in the first report. Their exact pedigree relationships were not assigned.

Sensitivity analyses retain the lexicographically first sample within each recorded family (1,260 representatives) or within each connected component of the union of the two reported kinship graphs (1,273 representatives). Selection does not use satellite measurements. The component rule is conservative and is not a maximum independent-set calculation. The reports are thresholded at first/second degree, so absence of a reported pair does not establish absence of more distant relatedness.

Sex comes from the corrected pedigree of the high-coverage release (1kGP.3202\_samples.pedigree\_info.txt) \cite{ByrskaBishop2022}. It differs from the 2013 pedigree in the sex of one sample: HG02300 is male in the 2013 file and female in the corrected one, and our data agree (X dosage 1.980, Y dosage 0.004). The two files also differ in 13 parent links; neither of the two affected parent--offspring links (HG02567--HG02569 and HG03466--HG03451) is in our pair set.

\subsubsection{Reference-panel representation}
We tabulated representation by target and Y lineage using the array-reference manifests, unique individuals in the ten-repeat CV predictions and the analysed cohort (Table~\ref{tab:reference-lineages}). Reference eligibility and read-based evaluation are counted separately; repeated predictions count once. Evaluated-individual counts differ by target (Table~\ref{tab:accuracy}). Grouping follows the archived ISOGG classification hierarchy, independently of the exploratory SNP-tree topology. Ten unique reference individuals lack calls in this lookup and remain unassigned; they are distinct from the two pending HGDP/SGDP calls in the geographic map. No call was inferred from a population label.

Representation is uneven at the displayed resolution. Among assigned references, O1a has a DYZ3 and DYZ19 reference but no Yq12 reference; O2b and several fine J groups have none for the three array targets. These counts identify limitations of the available panel, not the fraction of genetic diversity captured. Unassigned references limit statements about absent lineages. Leave-one-lineage-out evaluation tests prediction after excluding represented broad groups; it does not establish accuracy in every sparsely represented or unsampled subclade.

\begin{table}[p]
\caption{Y-lineage representation of the reference panels and analysed cohort. Each target column gives eligible assembly references / individuals evaluated in cross-validation. The Yq12 column also describes DYZ2; DYZ1, DYZ18S, YqHet2.7 and DYZ18sub add HG01261 (lineage R1b), whose Yq12 and DYZ2 lengths are excluded by the assembly-quality rule. Counts use the archived ISOGG 2019 classification hierarchy, independently of the inferred SNP tree. Unassigned denotes missing calls in the archived lookup, not a biological lineage. Reference individuals can also occur in the cohort; columns are not disjoint. A nonzero count does not establish adequate coverage of a lineage, and zero does not establish absence of a related reference in an ancestral or neighbouring clade.}\label{tab:reference-lineages}
\centering\small
\begin{tabular}{lrrrr}
\toprule
Y lineage & Cohort & DYZ3 & Yq12 & DYZ19 \\
\midrule
A0 & 3 & 1 / 1 & 1 / 1 & 2 / 2 \\
A1a & 5 & 3 / 3 & 3 / 3 & 3 / 3 \\
B & 8 & 3 / 3 & 2 / 2 & 3 / 3 \\
D & 20 & 5 / 2 & 1 / 1 & 2 / 2 \\
E & 466 & 24 / 20 & 16 / 16 & 32 / 27 \\
C & 32 & 3 / 1 & 2 / 1 & 6 / 3 \\
F4 & 1 & 1 / 1 & 1 / 1 & 1 / 1 \\
G & 28 & 3 / 3 & 3 / 3 & 4 / 4 \\
H & 71 & 4 / 3 & 2 / 2 & 4 / 3 \\
I & 62 & 4 / 3 & 3 / 2 & 3 / 3 \\
J1 & 19 & 3 / 3 & 2 / 2 & 3 / 3 \\
J2a1* & 1 & 0 / 0 & 0 / 0 & 0 / 0 \\
J2a1a1a & 17 & 2 / 1 & 1 / 1 & 1 / 1 \\
J2a1a1b & 25 & 3 / 2 & 3 / 2 & 3 / 2 \\
J2a1a2 & 3 & 0 / 0 & 0 / 0 & 0 / 0 \\
J2a2 & 1 & 0 / 0 & 0 / 0 & 0 / 0 \\
J2b & 24 & 1 / 1 & 1 / 1 & 1 / 1 \\
L & 31 & 1 / 1 & 1 / 1 & 2 / 2 \\
T & 11 & 2 / 2 & 1 / 1 & 2 / 2 \\
NO* & 1 & 1 / 1 & 0 / 0 & 1 / 1 \\
N & 32 & 5 / 5 & 5 / 5 & 6 / 6 \\
O1a & 20 & 1 / 1 & 0 / 0 & 1 / 1 \\
O1b & 87 & 4 / 4 & 2 / 2 & 4 / 4 \\
O2a & 139 & 14 / 13 & 11 / 11 & 14 / 14 \\
O2b & 2 & 0 / 0 & 0 / 0 & 0 / 0 \\
Q & 63 & 6 / 6 & 4 / 4 & 6 / 6 \\
R1a & 102 & 7 / 5 & 5 / 5 & 7 / 5 \\
R1b & 283 & 15 / 15 & 12 / 12 & 15 / 15 \\
R2 & 40 & 3 / 1 & 1 / 1 & 3 / 1 \\
Unassigned & 0 & 10 / 3 & 7 / 1 & 10 / 3 \\
\midrule
Total & 1597 & 129 / 104 & 90 / 81 & 139 / 118 \\
\botrule
\end{tabular}
\end{table}


\subsubsection{Geographic map and tree display}
Shared individuals are identified by exact sample identifiers after removing the SGDP manifest prefix and by documented SGDP individual aliases or sequencing identifiers. Eleven SGDP individuals also occur in 1kGP and eleven in HGDP, giving 2,302 unique individuals from 2,324 retained cohort records. Each counts once. Coordinates for shared individuals follow the fixed precedence 1kGP, HGDP, SGDP; discordant coordinates are not averaged. The overlap manifest retains all source identifiers, cohort memberships and coordinates.

Natural Earth 1:110 million land polygons provide the background. Identical reported coordinates rounded to six decimal places are pooled into 132 locations; nearby distinct locations remain separate. Radius is proportional to $\sqrt n$, making pie area proportional to location count. Overlapping pies are displaced with connectors to their reported coordinates. The map describes cohort coverage.

Targeted topology checks place HG03742 as sister to 280 N/O chromosomes, consistent with its published K2a assignment \cite{Poznik2016}, and recover monophyly of the 42 J2a1a1 calls. We checked sample coverage, the A0 root split, monophyly of the displayed rows, and preservation of unrooted splits and \vsette{total branch length (0.009555553) after rooting. Of the haplogroup labels, 267 of 309 Poznik and 1,452 of 1,485 HaploGrouper labels are monophyletic. Among included males, DYZ19, DYZ3 and Yq12 differed by more than 25\% in 4, 44 and 48 of the 545 sister pairs. The tree is not used} to infer state-transition counts or dates.

All five measurement columns use the same 1,597 read-based estimates. Reference status is not encoded by point appearance; the 136 males with reference assemblies remain identified in the sample metadata and sensitivity analyses. All four arm-imbalanced and six Yp-depleted profiles are excluded from population summaries\vsette{; the four males with a PAR2 gain leave only the Yq12 and ratio columns}. HG01967 and HG02966 are shown as diamonds only in Yq12 and DYZ19; ordinary points represent their other measurements; text labels identify substantial deviations using the display criteria below. Medians and interquartile ranges are computed independently within each row and target using only individuals included in population summaries. Only the DYZ3--R1b and DYZ19--J2a1a1 backgrounds are highlighted. Assembly spans appear in the assembly supplement and are not pooled with these population estimates.

\subsection{Supporting analyses}\label{sec:m-supporting}

\subsubsection{Details of the main results}\label{sec:res-main}
Among 123 completely assembled Y centromeres, the DYZ3 alpha-satellite array varied 13.8-fold in length, and among 84 complete Yq12 regions the heterochromatin varied from 10.3 to 53.1 Mb \cite{Hallast2026}. Earlier short-read studies measured DYZ3 in 366 males \cite{Miga2014}, DYZ1 in 396 \cite{Altemose2014} and DYZ1 coverage in 1,225 \cite{SanRoman2024}. The autosomal window for Y dosage (chr8:91--101 Mb) has keys taken from its projections into 101 high-quality HPRC haplotypes. Among 2,503 candidate windows ranked on pangenome conservation, single-copy state, copy-number stability, replicate stability, population spread and G+C match to the ruler, it was the best-ranked window that stayed within 1.89\% of the median of five windows in every sample, whereas the two higher-ranked windows, both on chromosome 11, shifted by up to 24.7\% in three males with a chromosome 11 gain.

The DYZ1:DYZ2 ratio had a median error of \vsette{2.61\%}. Short DYZ19 arrays were measured with \vsette{3.12\%} and \vsette{3.79\%} errors in two J2a1a1 males and \vsette{11.34\%} in the third. Parent--offspring signed medians were $-0.68\%$ for DYZ3, \vsette{$-0.09\%$} for Yq12 and $-0.45\%$ for DYZ19. In 243 females held out from training, the median background was 0.0128\% of the median reference DYZ3 array, \vsette{0.0202\%} of Yq12 and \vsette{0.0093\%} of DYZ19.

DYZ3 ranges were \vsette{226.2--532.6} kb in O2a2b1a1 and \vsette{754.8} kb--1.386 Mb in O2a2b1a2. In 518 further PCR-free HGDP and SGDP individuals with published calls \cite{Hallast2021}, DYZ3 medians were 364.0 kb in 50 R1b and 879.9 kb in 468 other males, and 265.4 kb in 13 O2a2b1a1 and 1.041 Mb in 14 O2a2b1a2 males. A YRI parent--offspring pair in E1a2a1b2 (NA18910 and NA18911) also carried short DYZ19 arrays (\vsette{83.2} and \vsette{77.4} kb). \vsette{The four PAR2 gains were HG02773, HG02895, HG03685 and NA11930. Independent libraries reproduced them with doses of 1.43--1.58 (Y state section), and the Yq12 of HG03685 was 1.14 times the lineage median. HG02773 had Yq12 of 56.37 Mb against 24.28 Mb in the parent HG02771, and NA11930 had 48.54 Mb against 21.96 Mb in the offspring NA10845, with DYZ3 within 1.1\%. Two further pairs differed in Yq12 by more than 25\%: NA12056--NA10851 by 37.1\% with DYZ3 within 0.6\% and DYZ19 within 3.0\%, and HG02317--HG02316 by 29.9\% with DYZ3 also 31.2\% apart. All four pairs cross sequencing batches, so library effects and changes during cell-line culture cannot be excluded.} Compared with its parent HG02965, HG02966 had near-absent Yq12 and DYZ19 signal and doses of \vsette{1.91--1.94} for X-degenerate blocks 1--5 and 0.0034--0.0039 for blocks 6--8, supporting an isodicentric Y that requires structural confirmation \cite{Porubsky2025}.

One annotated unit (one HMMER hit) spans a median of \vsette{3,478.9} bp of DYZ1 and \vsette{2,404.4} bp of DYZ2 in the 83 annotated assemblies \cite{Hallast2026}, so the base-pair ratio converts into a ratio of unit numbers: \vsette{0.967} in the median male (first to 99th percentile \vsette{0.678--1.462}). \vsette{The median of the 83 assemblies is 0.971 (mean 1.002), close to the published mean of 1.004 \cite{Hallast2026}.} The median unit ratio was \vsette{1.058} in haplogroup E (95\% interval \vsette{1.046--1.068}, resampling kinship components; \vsette{463} males), \vsette{12.5--17.2\%} above that of haplogroups D, I, C, O and R1b (20 to 282 males per haplogroup). The ratio error had a median of \vsette{2.63\%} per male and a 90th percentile of \vsette{7.65\%}, with no lineage bias: median signed ratio errors were \vsette{$-1.12\%$} in E (16 males), \vsette{$+0.11\%$} in O (13 males) and \vsette{$+0.02\%$} in R1b (12 males). Assembly medians were \vsette{1.046} for E, \vsette{0.950} for R1b and \vsette{0.923} for O (16, 12 and 13 assemblies). Unit lengths barely differ between lineages: across lineages with at least three of 90 annotated assemblies, the median base pairs per HMMER hit ranged from 3,470.8 to 3,491.3 for DYZ1 and from 2,400.9 to 2,409.8 for DYZ2, so lineage-specific unit lengths would change the unit ratio by at most 0.47\%, against lineage differences of \vsette{12.5--17.2\%}. DYZ19 below 100 kb rests on 4 assemblies (62.3--63.5 kb).

\subsubsection{Complete arrays in the Y pangenome define length and composition}\label{sec:res-truth}
\vsette{Complete arrays from the 142-individual Y pangenome \cite{Hallast2026} and HPRC release 2 \cite{HPRCRelease2} define the reference targets. After the target-specific quality checks, the panels contain 129 DYZ3, 90 Yq12, 91 DYZ1, 90 DYZ2, 91 DYZ18S, 139 DYZ19 and 140 DYZ17 targets. The corresponding evaluated sets contain 101, 80, 81, 80, 81, 115 and 115 individuals with short reads. Reference-panel size and the number of evaluated individuals are distinct (Table~\ref{tab:reference-lineages}).}

\vsette{For the 91 Yq12 annotations before the final region-specific exclusions, including 83 primary-resource arrays and eight HPRC replacements, the block definitions of Hallast et al. \cite{Hallast2026} assign median fractions of 57.13\% to DYZ1 and 40.64\% to DYZ2. DYZ1 spans 45.25--67.74\% of the reference region. Other annotated repeat classes and the unannotated remainder complete the partition; the residual unannotated fraction has a median of 0.57\% (161,070 bp). These are assigned-base fractions of the reference-coordinate span, not differences between independently inferred array boundaries. Extended Data Fig.~\ref{fig:yq12-composition} displays every annotated array.}

\vsette{Reference definitions distinguish spans from class amounts. DYZ3, DYZ19 and DYZ17 use array spans, whereas DYZ1, DYZ2 and DYZ18S use sums of retained annotation-block lengths. Total Yq12 is evaluated against the reference region span, and its prediction sums DYZ1, DYZ2 and a calibrated residual. Our HMMs define marker-sampling regions rather than replacement truth lengths. Calibration and evaluation use the published or equivalently extended annotation without marker-copy adjustment; the definitions and their development history remain conditions of the reported accuracy.}

\subsubsection{Yruler: marker k-mers and one pass over the reads}\label{sec:res-yruler}
\vsette{Yruler counts canonical 31-mers selected from complete arrays in the Y pangenome (Fig.~\ref{fig:method}). Class marker abundances are corrected with a G+C curve from single-copy Y markers and normalized to the same Y ruler. Training-array annotations supply the scales that convert those ratios to lengths. Independent markers of the chr8:91--101 Mb autosomal window supply Y dosage, with a factor of 0.9614 determined from reference males. Array length therefore measures repeat content relative to the ruler-bearing Y sequence, while dosage measures Y representation relative to the autosomal window.}

\vsette{The implementation reads BAM, CRAM, SAM, FASTQ and FASTA, including unmapped reads. Indexed BAM/CRAM supports parallel genomic shards; unindexed files and text formats are read sequentially, with parallel counting. For HG00126, an unindexed CRAM gave the same marker counts and estimates as the indexed file (930,909,168 primary records; 146.7 s with one reader and 84.0 s with the index, 16 threads); its 6,824,570 chrY and unmapped reads also gave identical results as BAM, interleaved FASTQ or paired FASTQ files. Counting depends on read sequence, without requiring unique placement inside an array. Saved per-marker counts support recalibration without rereading the genome file.}

\vsette{The released model uses a masked-superstring marker archive \cite{Sladky2023} that occupies 3.97 MB (3,974,793 bytes) and represents 8,229,748 marker memberships: seven reported satellite targets, the Yq12 residual helper, the Y ruler and the chr8-window markers. A k-mer can belong to more than one class, so memberships do not count unique sequences across classes. The release parameter file records each class's scale and the Yq12 component-sum expression, the dosage scale of 0.9614 and the no-Y cut of 0.0067. The cohort estimates use this archive and parameter file; evaluation refits marker sets and scales within training partitions while keeping the target definitions and output expression fixed.}

\subsubsection{Prediction accuracy and generalization}\label{sec:res-benchmark}
Prediction accuracy was assessed using ten repeats of family-grouped five-fold CV, with markers, presence and outside thresholds and scale fitted inside training folds (Fig.~\ref{fig:method}b). Pooled median absolute errors were \vsette{2.32\%} for DYZ3, \vsette{1.72\%} for Yq12, \vsette{2.48\%} for DYZ19. Evaluated sample counts were \vsette{101, 80 and 115}, respectively. The derived DYZ1:DYZ2 ratio had a pooled median error of \vsette{2.61\%} in \vsette{80} individuals. Excluding entire lineages gave pooled median absolute errors of \vsette{2.39\%, 1.89\% and 2.63\%}, respectively (Extended Data Fig.~\ref{fig:lineage-validation}). At \vsette{1.0$\times$} median autosomal coverage, pooled medians were \vsette{3.28\%, 2.06\% and 4.21\%} (Extended Data Fig.~\ref{fig:s-depth}). Estimates of the 115 calibration males in the population analysis remain in-sample; those of the 21 held-out males with reference assemblies do not.

The Miga/Altemose comparison informed method design \cite{Miga2014,Altemose2014} and is reported as development evidence. \vsette{Twenty-one males whose reads were never used in any choice gave 1.55\% for DYZ3 (18 males), 0.51\% for Yq12 (3) and 2.90\% for DYZ19 (14).} Independent-cohort performance was not assessed.


\subsubsection{Related satellite classes differ in measurability}\label{sec:res-classes}
Grouped CV errors are class-specific: \vsette{DYZ18S has a pooled median error of 3.92\% (Table~\ref{tab:accuracy}) and a largest single-prediction error of 17.63\%; its mean signed error is $-14.05\%$ to $-16.75\%$ in three of the four haplogroup A references and $+0.82\%$ in the fourth (HG02984).} These tails remain limitations despite the low median errors of the main arrays.

\vsette{YqHet2.7 and DYZ18sub are not reported: on the published annotation \cite{Hallast2026} their CV median errors were 13.11\% and 5.63\%, with maxima of 313.69\% and 33.52\%, because the YqHet2.7 motif hits fall mostly inside DYZ1 and the class boundary depends on a score comparison that 31-mers cannot reproduce.} The maximum DYZ18S error remains \vsette{17.63\%} in HG02666. \vsette{With an earlier model,} HG002 had a mean signed DYZ2 error of $-1.81\%$ with the G+C floor and $-12.89\%$ without it, consistent with the loss of G+C-poor reads in its library. No single summary establishes equal reliability across all targets or sequencing libraries.

Reference annotation can partition Yq12 into assigned sequence and unassigned residuals. \vsette{Read estimates of DYZ1, DYZ2 and DYZ18S are separate measurements.} Validation at the short or near-absent state remains particularly important: estimates far below the training range identify signal depletion for follow-up, not base-pair-resolved deletion boundaries.

\subsubsection{Aligners, platforms and read groups}\label{sec:res-groups}
On 115 individuals with both NYGC BWA CRAMs and DRAGEN 3.7.6 BAMs, whole-file counting with the development marker set produced identical summary files. Because all eligible read sequences were counted, differences in where those reads aligned did not affect the estimates. \vsette{Aligned read density did change (Extended Data Fig.~\ref{fig:aligner}): in 130 males with both alignments, density in GRCh38 array blocks relative to chrY:14.5--15.5 Mb was lower with DRAGEN by medians of 12.80\% for DYZ1, 5.23\% for DYZ19, 3.53\% for DYZ2 and 0.37\% for DYZ3, and after removing one offset per array the differences between individuals had medians of 1.80\%, 1.22\%, 0.27\% and 0.14\% and maxima of 9.91\%, 6.86\%, 2.03\% and 1.37\%.} This comparison concerns two representations of the same sequencing data; it does not test independent libraries or sequencing instruments.

For indexed DRAGEN BAMs, we also tested a reduced set of genomic regions that captures reads carrying the array and depth markers, together with unmapped reads. We chose additional windows on 30 individuals and evaluated the resulting BED on a further 100. Across the eight satellite classes of this development test, the range of the BED-to-whole normalized abundance ratio was 0.26--0.42 percentage points over all 130 individuals. Median ratios were 1.0018--1.0055, so a separate calibration is needed for this mode. \vsette{With the final markers, the one sample tested with a DRAGEN BAM and the BED (HG00126) differed from its whole-CRAM estimates by $-0.021\%$ to $+0.214\%$ across the seven classes, and the whole CRAM matched the cohort table exactly; a NYGC CRAM read through this BED gives DYZ2 85.84\% too low, so the BED applies only to DRAGEN 3.7.6 alignments.} Mean processing time was 22.32 s for the restricted regions, compared with a recorded whole-file mean of 124.04 s. These are local-file measurements and do not include a cohort download.

In a separate 16-thread whole-CRAM optimization experiment on HG00126, the ten-class prototype processed the file in 79--86 s, with peak resident memory of 5.4--6.2 GB depending on buffering. This establishes a computational baseline rather than a final cohort-wide runtime claim.



\draftnote{Repeat the alignment and restricted-region checks with the final markers and mode-specific calibration. Final runtime reporting needs a fixed software build, hardware description and separation of compute time from remote I/O. Long-read and Hi-C results require their own validated calibrations before quantitative claims are added.}

\subsubsection{Y dosage and sample quality}\label{sec:res-dosage}
The 41 female samples had Y dosage at most \vsette{0.0053}. Of 1,598 male samples, 1,597 had estimates; NA19226 had none (one X, no Y signal, Y lost in the sample of the parent of an XY offspring). Low-dose samples were retained when their estimates were available; arm-imbalanced, Yq-loss and Yp-depleted profiles were excluded from population summaries. X dosage was \vsette{0.977--1.021} in 100 reference XY males and had a median of \vsette{1.958} in the 41 females; it identified one XXY male (HG01683: X \vsette{1.983}, Y \vsette{0.987}), which the Y-based rule alone had not flagged, and confirmed one X in the two males with Y dosage near 2 (HG02869 and HG03456: X \vsette{0.991} and \vsette{1.010}). X dosage in the \vsette{1,561} XY males with Y state normal was \vsette{0.965--1.027}. Pseudoautosomal dosage counts the X and Y copies together (1 for two copies): it was \vsette{0.938--1.068} for PAR1 and \vsette{0.956--1.042} for PAR2 in the same 100 reference males, \vsette{1.466--1.526 and 1.397--1.511} in HG01683, HG02869 and HG03456 (three copies) and \vsette{0.507 and 0.513} in NA19226 (one copy). Among the XY males, seven had a PAR2 dosage consistent with one copy more or fewer (a partial change is not excluded) (HG01589, HG03809 and HG04149 at \vsette{0.514--0.616}, besides HG02966 at 0.513, which is counted with the arm-imbalanced samples; HG02773, HG02895, HG03685 and NA11930 at \vsette{1.438--1.471}) and NA20864 had a PAR1 dosage of \vsette{1.230}.

Normalizing satellite counts to single-copy Y abundance separates estimated length per Y from the amount of Y-bearing material, but cannot rescue negligible Y coverage. Dosage below 0.9 is a descriptive flag rather than an exclusion criterion. Elevated dosage also motivates inspection: HG01967 has dosage \vsette{1.3150} and DYZ19 \vsette{87,370} bp and is retained as an unusual estimate for follow-up, outside population summaries. Neither low nor high dosage is assigned a cause from this ratio alone. In Fig.~\ref{fig:biology}a, satellite-length and length-ratio labels identify up to three observations per pooled 1.5-interquartile-range tail, additionally requiring at least 50\% deviation from the median of the displayed clade. HG01967 and HG02966 are also labelled in Yq12 and DYZ19 when they meet this minimum deviation; the parent--offspring pair NA18910--NA18911 is similarly labelled in DYZ19. Dosage labels identify all observations above 1.30. These display thresholds suppress labels for small deviations in narrow distributions; they are descriptive choices, not validated biological or technical cutoffs, and do not establish measurement error or account for pedigree dependence.

\subsubsection{Checks without assembly lengths}\label{sec:res-truthfree}
In \vsette{243} female genomes held out from marker selection (\vsette{2,430} predictions), the false signal of each class had medians of \vsette{30 bp (DYZ19) to 6,969 bp (DYZ2) and maxima of 1,431 bp (DYZ18S) to 127,935 bp (Yq12)}. Every 1kGP male except HG02966 lay at least \vsette{20.9} times above the female maximum in every class; HG02966 lay at or below it in \vsette{five} classes. \vsette{Mixing reads of six male--female pairs at Y fractions from 1 to 0.01 and base coverage of 5, 1 and 0.2$\times$ (three seeds, 270 runs) kept the ruler basis in every run, with a lowest dosage of 0.0082 above the no-Y cut (0.0067), and no run of females alone used it (54 runs, highest dosage 0.0044). Median dosage followed the fraction (0.1015--0.1028 at 0.1, 0.0123--0.0138 at 0.01). Lengths per Y changed, relative to the whole genome of the same male, by medians across the seven classes of 0.23--1.93\% at 5$\times$, 0.88--2.61\% at 1$\times$ and 2.04--8.55\% at 0.2$\times$ at fraction 1; 1.17--5.18\%, 2.58--12.72\% and 3.49--22.09\% at fraction 0.1; and 6.94--17.33\%, 10.64--31.60\% and 25.02--98.71\% at fraction 0.01. The G+C curve was lost at 0.2$\times$ at every fraction and at 1$\times$ from fraction 0.1. Y signal is therefore detected down to a fraction of 0.01, but length errors depend on the coverage of Y-bearing reads: at 0.5$\times$ (fraction 0.5 at 1$\times$ or 0.1 at 5$\times$) medians were at most 6.83\%.} In 114 reference males, thinning marker counts by read G+C bin, each count on its own (keep probabilities 0.2, 0.4, 0.6, 0.8 and 0.9 below 20\%, 20--25\%, 25--30\%, 30--35\% and 35--40\% G+C) changed the estimates of all males by one factor per class (90th percentile minus median at most \vsette{0.46} percentage points). The median change was \vsette{2.41--14.42\%} without G+C correction and \vsette{0.73--7.46\%} with it; \vsette{Yq12 fell from 7.03\% to 0.73\%} and DYZ19 from \vsette{14.42\% to 7.46\%}.

\subsubsection{Library preparation and sequencing batches}\label{sec:res-libraries}
\vsette{Ten} males of eight populations sequenced by both SGDP (PCR-free, 100 bp reads) and the 1kGP high-coverage project (PCR-free, 150 bp reads) differed by medians of \vsette{1.89--5.26\% in six classes} and by \vsette{9.80\%} in DYZ19\vsette{, higher in SGDP in all ten}. \vsette{Cutting the 1kGP reads of these males to 100 bases raised DYZ19 by a median of 5.19\% (3.84--7.79\%), DYZ18S by 3.52\% and the other classes by at most 1.50\% in absolute value, and against the cut reads the SGDP excess in DYZ19 fell from 9.80\% to 4.71\%.} \vsette{Counting the overlapping bases of read mates once on both sides left the DYZ19 difference at 10.48\% (higher in SGDP in all ten pairs) and moved the other six classes by at most 0.18 points, so mate overlap does not explain it.} \vsette{Eleven} males sequenced by SGDP and by HGDP, all Papuan, differed by \vsette{55.96\%} in DYZ2. These HGDP libraries were PCR-amplified: 151 of the 552 HGDP males have PCR libraries, each sequenced on one lane, and the others PCR-free libraries multiplexed over lanes \cite{Bergstrom2020}. In \vsette{13} populations with at least three males of each kind, PCR and PCR-free population medians differed by \vsette{$-45.5\%$} for DYZ2 (negative in all \vsette{13}), \vsette{$+24.7\%$} for DYZ18S and \vsette{$-12.7\%$} for Yq12. SGDP sequencing panel B (22 samples) used a library preparation other than the PCR-free protocol of panel C \cite{Mallick2016}. Population summaries therefore use PCR-free libraries. A drop of coverage at low G+C (ruler curve below 0.975 at 8 of 31 G+C bases) marked 81 of 87 PCR and 6 of 236 PCR-free libraries among the 323 HGDP males with a G+C curve, whereas the duplicate fraction was lower in PCR libraries (median 6.5\%) than in PCR-free libraries (15.1\%), although all were sequenced on HiSeq X Ten. All 1kGP libraries are PCR-free and were sequenced on 133 NovaSeq flowcells. \vsette{In 282 cross-batch parent--offspring pairs, differences followed the difference in mate overlap (median overlap share 0.061, Spearman $\rho=+0.292$ for DYZ19, $+0.288$ for DYZ18S, $+0.236$ for DYZ17 and $+0.210$ for DYZ3). Counting overlapping bases once, which removed a median of 0.32\% of the bases, moved CV medians by at most 0.09 points (DYZ19 2.48\% to 2.57\%) and left these correlations unchanged in 60 pairs, so double counting does not explain them.} \vsette{Yruler counts every read: in 30 males, removing reads flagged as duplicates moved estimates by medians of $-16.18\%$ (DYZ2) to $+7.89\%$ (DYZ3), and keeping only reads with mapping quality of at least 20 removed 27.82\% (DYZ18S) to 99.39\% (DYZ3) of the signal, which shows why these filters are not applied to satellite reads.} In \vsette{297} parent--offspring pairs, the parent--offspring difference followed the difference of their low-G+C curve values for DYZ18S (Spearman $\rho=\vsette{-0.509}$), DYZ1 (\vsette{$-0.393$}), DYZ17 (\vsette{$-0.382$}) and DYZ2 (\vsette{$+0.377$}), but not for DYZ3 (\vsette{$0.000$}) or DYZ19 (\vsette{$-0.126$}): part of the error in these \vsette{four} classes is library G+C bias that the 31-mer G+C curve leaves. \vsette{In 81 CV males, the error shared by DYZ1, DYZ17 and DYZ18S (standard deviation 3.98\%) followed the low-G+C curve (Spearman $\rho=-0.523$) and not the fraction of absent markers ($-0.132$ for the ruler, $-0.034$ to $+0.008$ for these classes).} \vsette{Yq12, DYZ1 and DYZ2 also move with X dosage (CV slope 1.59 for Yq12; parent--offspring slope 1.25, 95\% interval 0.83--1.57, in 270 cross-batch pairs with a normal state, absolute log differences below 0.25 and X dosage differences below 0.1, which leaves out the XXY male HG01683 and the PAR2-gain pairs; 0.099 in all 282 cross-batch pairs). A linear correction, fitted on training males, lowered the held-out Yq12 median from 1.72\% to 1.45\% but worsened parent--offspring agreement (1.82\% to 1.87\%) and the SGDP--1kGP pairs (3.71\% to 4.23\%), and failed the rule written before the test, so it is not applied.} A correction by the G+C of the first 100 bases of each read was not better than the 31-mer curve and did not remove the PCR effect on DYZ2.

\subsubsection{Pedigree agreement and empirical relatedness}\label{sec:res-pedigree}
The 309 parent--offspring pairs had median absolute relative differences of \vsette{1.86\%} for DYZ3 and \vsette{1.81\%} for Yq12. Signed medians were \vsette{$-0.68\%$ and $-0.09\%$}, respectively. The 282 original-parent/additional-offspring pairs had signed medians $-0.83\%$ and \vsette{$-0.05\%$}; the 25 same-batch pairs had \vsette{$+0.22\%$ and $+0.29\%$}. Median absolute differences were \vsette{2.85\%} for DYZ19, \vsette{2.96\%} for DYZ1, \vsette{2.61\%} for DYZ2 and \vsette{4.44\%} for the DYZ1:DYZ2 ratio (Extended Data Fig.~\ref{fig:all-pedigrees}). Seventy-eight pairs involve a reference individual. One factor per pair, shared by the seven classes, explained 57.9\% of the variance of the log differences in 273 cross-batch pairs (standard deviation 3.89\%) and 51.0\% in 22 same-batch pairs (2.81\%), so a comparison of two individuals carries a shared library error that the per-class widths do not describe. Pedigree agreement combines measurement error and biological variation; it is not an independent measure of accuracy against assemblies.

Public KING estimates support all \vsette{309} recorded parent--offspring relationships. Ten additional close-kinship links across different FamilyIDs occur in both reports, with two further links present in only one report. Retaining one included representative per recorded family gives 1,260 samples; using connected components of both kinship reports gives 1,273. These sensitivity subsets retain the DYZ19 separation and similar DYZ1:DYZ2 summaries (Table~\ref{tab:relatedness}).

\vsette{More distant kinship came from the HGDP and 1kGP relatedness of the harmonized gnomAD resource \cite{Koenig2024} (pairs with kinship at least 0.050) and, for the 74 SGDP individuals outside it, from KING-robust kinship on 129,632 random autosomal SNPs of the published SGDP genotypes \cite{Mallick2016}. In 1kGP it added 11 third-degree links and no closer one; one male per enlarged component (1,260 males) left the DYZ1:DYZ2 standard-deviation ratio at 0.762, the Hodges--Lehmann ratios at 0.340 (DYZ3, R1b against other males) and 0.235 (DYZ19, J2a1a1 against other J males) and $H$ for DYZ3 at 791.0 (794.8 before). Among the 518 HGDP and SGDP individuals, 11 pairs were related at third degree or closer (three second-degree, none first-degree); one per component (509 individuals) gave standard-deviation ratios of 0.801 in HGDP (0.798 before) and 0.655 in SGDP (unchanged), 22 short DYZ19 arrays among 62 J chromosomes, all J2a1a1, and a DYZ3 median of 360.6 kb in 49 R1b males (364.0 kb in 50 before). Close paternal relatives can also share an almost identical Y without autosomal kinship. Grouping the cross-validation folds by distance on the Y tree as well as by family (single linkage below $2\times10^{-5}$ substitutions per site, which gave 69 groups for the 115 males; rule fixed before the run) changed median errors by $-0.08$ to $+0.11$ points (DYZ3 2.36\%, Yq12 1.75\%, DYZ1 2.29\%, DYZ2 2.42\%, DYZ18S 4.02\%, DYZ19 2.42\%, DYZ17 2.35\%) and 90th percentiles by $-0.19$ to $+0.16$ points, so near relatives in training do not inflate the reported accuracy. Only two pairs of male siblings were available, too few for a separate concordance estimate.}

HG01527 and HG01529 are parent and offspring (IBS010), supported by KING kinship of 0.253 and 0.250 and IBS0 below 0.00004. Their DYZ18S estimates are \vsette{590,804 and 584,440} bp and provide one familial observation. HG02965 and HG02966 are also a supported parent--offspring pair (NG20). The reference parent has Yq12 \vsette{27.54} Mb and DYZ19 \vsette{238.4} kb; the offspring has \vsette{38.8} kb and \vsette{266} bp. These values indicate marker depletion far below the calibration range; structural boundaries require sequence confirmation.

The NG20 trio comprises parent HG02965, parent HG02964 and offspring HG02966. DYZ3 estimates are 0.951 Mb and 1.560 Mb in parent and offspring, with Y dosage \vsette{0.9965 and 1.1862}. Multiplying $\widehat L_{ic}$ by $D_i$ cancels their shared Y denominator. The offspring-to-parent autosomal-normalized DYZ3 signal ratio is \vsette{1.95}; exact doubling would predict an offspring length of \vsette{1.597} Mb. Near doubling together with distal satellite depletion motivates an isodicentric Y hypothesis \cite{Lange2009}, but reuses the same marker data. Junction-spanning sequence and a separate DNA source are needed to establish chromosome structure and distinguish donor variation from cell-line change. HG01967 is parent of XX offspring HG01969 in a different trio (parent HG01968), which supplies no paternal Y transmission comparison.

\begin{table}[tbp]
\caption{Relatedness and reference-panel sensitivity. J-short counts J2a1a1 calls and J-other the remaining J calls; their DYZ19 ranges are separated in every subset. Ratio is the median DYZ1:DYZ2 base-pair ratio, and $\rho$ is the Spearman correlation. \vsette{Yq12 and ratio columns here include the four males with a PAR2 gain.} Representatives were selected lexicographically without using satellite measurements. KING components use the union of two public first/second-degree reports and do not establish absence of distant kinship.}\label{tab:relatedness}
\centering
\small
\begin{tabular}{lrrrrr}
\toprule
Subset & $n$ & J-short & J-other & Ratio & $\rho$ \\
\midrule
Population analysis subset & 1,587 & 42 & 44 & \vsette{1.399} & \vsette{0.818} \\
Outside reference panel & 1,451 & 37 & 40 & \vsette{1.398} & \vsette{0.819} \\
Y dosage at least 0.9 & \vsette{1,573} & 41 & 44 & \vsette{1.399} & \vsette{0.817} \\
One per recorded family & 1,260 & 37 & 37 & \vsette{1.394} & \vsette{0.829} \\
One per KING component & 1,273 & 37 & 37 & \vsette{1.394} & \vsette{0.827} \\
\botrule
\end{tabular}
\end{table}

Transmission agreement combines technical effects and possible biological variation. The few-kilobase DYZ2 variants resolved in CEPH 1463 \cite{Porubsky2025} illustrate changes below the percent-level resolution of a typical Yq12 measurement. Neither pedigree agreement nor an apparent directional difference establishes a mutation rate.

\subsubsection{Array lengths in 1000 Genomes males}\label{sec:res-catalog}
All 1,597 male-labelled samples with estimates have \vsette{seven} satellite measurements, with 136 males with reference assemblies (115 used for calibration, 21 held out) and 1,461 other individuals (Fig.~\ref{fig:biology}; Extended Data Fig.~\ref{fig:s-population}). Population summaries exclude ten arm-imbalanced or Yp-depleted samples and retain concordantly low-dose profiles with available estimates. Median lengths in the 1,587 included individuals were \vsette{857,643} bp for DYZ3, \vsette{332,356} bp for DYZ19 \vsette{and 861,548 bp for DYZ17, and in the 1,583 males of the Yq12 summaries 26,924,721} bp for Yq12, \vsette{15,630,668} bp for DYZ1 and \vsette{10,740,802} bp for DYZ2. These values describe the analysed sample, whose population representation is uneven.

DYZ3 had a median of \vsette{318.8 kb} in 282 R1b individuals and \vsette{980.9 kb} in the other 1,305, recovering the reported association \cite{Hallast2023}. H had the largest broad-lineage median (\vsette{1.349} Mb; 71 individuals); R1a had \vsette{540,670} bp in 100 individuals. O had median Yq12 \vsette{34.7} Mb. The 465 included E chromosomes had DYZ3 and DYZ19 medians of \vsette{1,091,996 and 242,090} bp, respectively.

HG04006 has the largest DYZ3 estimate in the analysed sample, \vsette{2.768} Mb, beyond the longest calibration array; independent measurement is needed. Parent HG02684 and offspring HG02686 have DYZ3 estimates of 2.51 and 2.58 Mb, with DYZ1:DYZ2 base-pair ratios \vsette{2.44 and 2.62}. Their resemblance is one familial observation.

\subsubsection{Lineage tests and correlation between arrays}\label{sec:res-tests}
\vsette{In 1,271 males, one per kinship component without the four males with a PAR2 gain, Kruskal--Wallis tests across the 14 lineages with at least 20 males gave $P<10^{-44}$ for all seven arrays and the unit ratio ($H$ 794.8 for DYZ3, 999.3 for DYZ19, 312.3 for Yq12, 247.2 for the ratio); Dunn tests with Holm correction left 25 (ratio) to 59 (DYZ19) of 91 lineage pairs below 0.05 (results table dunn\_pairs.tsv). These tests treat males of a lineage as independent, so they are descriptive; the phylogenetic models below address shared ancestry. The 41 males above the longest assembled DYZ3 array are mostly E (25) and H (9); the DYZ3 panel holds three H males, all at most 1.573 Mb. Without these 41 males, $H$ for DYZ3 in all males with lineages fell from 982.0 (1,587 males) to 951.9 (1,546 males), and no lineage median moved by more than 26 kb. Aligned read density in chrY:10,083,500--10,621,500 of GRCh38 over that in chrY:14.5--15.5 Mb, which uses no Yruler marker, followed the DYZ3 estimates of 1,587 males (Spearman $\rho=0.9992$, 0.9805 in the 41 males). A line fitted below 1,673,237 bp predicted these 41 males with a median difference of 0.58\% and a maximum of 2.85\% (HG04006: 2.732 against 2.768 Mb). This shows that read abundance stays proportional above the assembled range; it does not measure a physical array. The same density in GRCh38 blocks followed the estimates of Yq12, DYZ1 and DYZ2 in 1,583 males (Spearman $\rho=0.924$, $0.972$ and $0.813$); the saved bins have no DYZ18S or DYZ17 hits, so these two classes were not checked. The Hodges--Lehmann ratio of R1b to other chromosomes was 0.339 for DYZ3 (bootstrap 95\% interval 0.329--0.351), and that of J2a1a1 to other J was 0.235 for DYZ19 (0.230--0.241). Arrays vary almost independently: within lineages, Spearman correlations of DYZ3 with Yq12, DYZ1, DYZ2, DYZ19 and DYZ17 were $-0.004$, $+0.014$, $-0.018$, $+0.082$ and $-0.052$. The largest values outside the Yq12 components were $+0.327$ for DYZ18S with DYZ17 and $+0.256$ for DYZ1 with DYZ17. Shared error inflates both: CV errors of these pairs correlate at $+0.68$ and $+0.88$, and in 75 assemblies the within-lineage values are $+0.152$ and $+0.339$. Across individuals DYZ3 and DYZ19 correlate at $-0.463$; this disappears within lineages and across the medians of 14 lineages ($-0.152$, $P=0.60$; $-0.021$ without E and R1b), so it reflects the sizes of those two lineages. DYZ17 ranged from 577.2 kb to 1.358 Mb, with lineage medians from 787 kb in J to 953 kb in N.}

\vsette{On the tree, in the same 1,271 males, shared-branch covariance was strong for every array: maximum-likelihood $\lambda$ was 0.99 for DYZ3, 0.97 for Yq12, DYZ1, DYZ2 and the ratio, 0.96 for DYZ18S, 0.93 for DYZ17 and at the upper bound (0.9999) for DYZ19. Measurement error is not separated from biological residual variance, and the tree is not dated. Treating DYZ19 below 100, 125 or 150 kb as a state, parsimony needs at least two changes on the tree, consistent with J2a1a1 and the E1a2a1b2 pair.}

\subsubsection{Population tests of compositional balance}\label{sec:res-balance}
DYZ1 and DYZ2 estimates across the 1,583 population-summary individuals are correlated (Spearman $\rho=\vsette{0.8165}$), with median base-pair ratio \vsette{1.3990} and first--99th percentiles \vsette{0.9816--2.1151}. Family and public-kinship pruning, outside-reference-panel estimates and the earlier dosage threshold give similar ratio summaries (Table~\ref{tab:relatedness}).

The 91 assembly annotations provide a separate composition view: DYZ1 base-pair fractions range from \vsette{45.2\% to 67.7\%}, and the assigned classes plus unannotated sequence partition each annotated span (Extended Data Fig.~\ref{fig:yq12-composition}). Those fractions do not directly establish equal repeat-unit counts. Balanced DYZ1:DYZ2 unit abundance and evolutionary coupling were already reported in the assembly resource \cite{Hallast2026}.

\vsette{Within broad lineages, in 1,196 males without a reference assembly, one per kinship component, the standard deviation of log length was 0.194 for DYZ1 and 0.255 for DYZ2 (ratio 0.762, bootstrap 95\% interval 0.724--0.802; Pitman--Morgan test of equal variances $P=5\times10^{-67}$). The ratio was 0.798 (0.750--0.845) in 382 HGDP males and 0.655 (0.589--0.741) in 136 SGDP males, both PCR-free, with lineages from the first letter of published calls \cite{Hallast2021} (0.788 and 0.669 with the first three characters of the calls), and 0.714 (0.634--0.829) in the annotated DYZ1 and DYZ2 amounts of 75 assemblies \cite{Hallast2026} (0.716, 0.635--0.833, with annotated unit counts), without HG01261, whose Yq12 and DYZ2 targets were excluded. CV standard deviations of the log errors (0.033 and 0.039) are too small to produce this difference. The slope of the log ratio on the mean of the two log lengths is zero when the two variances are equal and shares no term with the rest of Yq12. Generalized least squares on the tree with Pagel $\lambda$ estimated by maximum likelihood gave $-0.309$ (95\% interval $-0.338$ to $-0.281$; $\lambda=0.974$; $P=1.6\times10^{-86}$) in the 1,196 males and $-0.330$ ($-0.436$ to $-0.224$; $P=2.9\times10^{-8}$) in the 75 assemblies; on log DYZ3 the slope was $+0.006$ ($-0.025$ to $0.037$; $P=0.69$), and on the log of the rest of Yq12 $+0.276$ ($0.185$ to $0.367$), so the decline does not come from the rest of Yq12. Slopes on log Yq12 ($-0.276$ within lineages and $-0.300$ on the tree in these males, $-0.330$ on the tree in the assemblies) are not used as evidence, because Yq12 contains both arrays: permuting DYZ2 among males of a lineage, which removes its correlation with DYZ1, gave a median slope of $-0.172$ (central 95\% $-0.201$ to $-0.139$; 1,000 permutations). Offspring were no closer to unit balance than parents (offspring closer in 159 of 297 pairs; Wilcoxon $P=0.72$; cross-batch pairs 145 of 273).} The remaining population test is whether supported departures from a declared unit-equivalent balance exceed joint annotation and measurement uncertainty and recur in independent data. It must address batch, cohort and family dependence. The present correlation or a ratio above one is insufficient to claim a novel exception, functional constraint or selection.

\subsubsection{Short DYZ19 arrays in sampled J2a1a1 chromosomes}\label{sec:res-lineages}
\vsette{A mixture of log DYZ19 lengths supported more than one state (BIC 1,086 for one component, 43.5 for two and $-479$ for three, with means of 65.8, 255.4 and 379.8 kb), and no estimate fell between 83.2 and 161.0 kb. The four B3 chromosomes (two B3a2, two B3b, from GWD and MSL) had 161.0--174.4 kb, among them HG03225, whose clean assembly array is 165.3 kb, against 284.4--299.4 kb in the four B2 chromosomes.} All 42 sampled J2a1a1 chromosomes had DYZ19 estimates of \vsette{59,169--74,656} bp, whereas the other 44 included J chromosomes spanned \vsette{220,914--305,264} bp. The short-state calls split between J2a1a1a (17 individuals) and J2a1a1b (25). Five parent--offspring pairs carried the short state in both members (HG01672--HG01671, HG01781--HG01780, HG02490--HG02492, HG02733--HG02735 and HG03905--HG03906; 61.3--74.7 kb, differences $-6.8\%$ to $+7.0\%$). Thirty-seven short-state individuals have no reference assembly, and 39 are outside the 115 calibration males. \vsette{Without Yruler markers, aligned read density per base in chrY:20,054,500--20,351,000 of GRCh38 over that in chrY:14.5--15.5 Mb was 0.111--0.144 in the J2a1a1 males and 0.645--0.896 in the other J males, 0.243 and 0.240 in NA18910 and NA18911, and correlated with the estimates in 238 males (42 J2a1a1, 44 other J and 152 others; Spearman $\rho=0.990$). In the whole GRCh38 block (chrY:20,054,500--21,184,000, located with marker hits in ten males), which includes flanking sequence, the values were 0.759--0.809 and 0.913--0.999 ($\rho=0.972$). Alignment to GRCh38 shows that the region is reduced; it gives neither a length nor a breakpoint. HaploGrouper with the ISOGG 2019 tree placed all 22 short HGDP and SGDP males in J2a1a1 (10 J2a1a1a and 12 J2a1a1b) and none of the 41 other J males, and matched the published calls of 12 1kGP controls.} Retaining one per documented family or public-kinship component leaves 37 short-state versus 37 other J chromosomes with the same separation.

This extends the short DYZ19 observation in HG02492 \cite{Hallast2023} to a larger sampled set. Four quality-passing J2a assembly references, HG02492, HG02735, NA20805 and NA20809, have spans of \vsette{62,251--63,499} bp, compared with a median of \vsette{321,716} bp across the 139 DYZ19 references; the J2b reference NA20905 measures \vsette{267,723} bp. Reference-panel observations and population evidence are reported separately.

NA18910 and NA18911, a parent--offspring pair from the YRI cohort (E1a2a1b2), also have short DYZ19 estimates of \vsette{83,190 and 77,375} bp, with Y dosage \vsette{1.0199 and 1.0112}. Public KING reports support their relationship. They provide one familial observation outside J2a1a1; the current sample contains no other E1a2a1b2 individuals. Measurements in separate sequencing batches do not establish a short state throughout that lineage. \vsette{Published haplogroups from Hallast et al. \cite{Hallast2021} were matched to 640 of the 665 HGDP and SGDP profiles. After keeping PCR-free libraries and estimable balanced profiles and removing 1kGP aliases and HGDP--SGDP duplicates, 518 individuals remained. Of their 63 J chromosomes, 22 had DYZ19 estimates of 65.5--74.9 kb and the others at least 250.8 kb; no chromosome of another lineage fell below 100 kb. All 22 carry labels beginning J2a1a in that ISOGG version; two males labelled J2a1a and two labelled J2a1a2b2a2a2 did not have short arrays, and these labels are not harmonized with the classification used for 1kGP.}

The \vsette{SNP tree} recovers a clade containing all 42 J2a1a1 chromosomes; J2a1a2 lies outside that clade (Extended Data Fig.~\ref{fig:j-dyz19}). This supports a lineage-associated short-array state in the sampled J2a1a1 chromosomes. The separate short-state YRI family and remaining fine-scale classification discordances require additional checks before inferring origins or transitions. The current tree is not time-calibrated.

\subsection{Interpretation of the supporting analyses}
The Y pangenome makes satellite variation accessible as assembled sequence across individuals, but an assembly for every sequenced genome remains impractical. Yruler uses complete arrays from this pangenome to define measurement targets and select markers, then estimates array length directly from sequencing-read abundance. The implementation measures total DYZ3 and Yq12 abundance alongside satellite sequence components. Measuring several targets in the same pass also exposes differences among arrays that would be obscured by a single estimate of Y heterochromatin content. The pangenome supplies sequence diversity and calibrated references. Grouped cross-validation assesses prediction within the current reference collection; lineage exclusion tests transfer to represented broad groups; performance in independent cohorts remains to be established.

The main contribution is the combination of multi-array measurement, explicit assembly calibration, cross-validated evaluation and Y dosage. Earlier studies already measured DYZ3 or DYZ1 from short reads \cite{Miga2014,Altemose2014}, and DYZ1 coverage has been studied in more than one thousand males \cite{SanRoman2024}. An empirical comparison must therefore assess accuracy and generalization on shared individuals, while acknowledging differences in sequencing depth and library preparation. Comparisons of previously published estimates with contemporary reads are not controlled comparisons of algorithms alone. Likewise, selecting a similar centromere reference haplotype \cite{Shiraishi2026} addresses a different target from directly measuring its length.

Our reference annotation emphasizes that an array does not have a unique length until its boundaries and sequence classes have been specified. A reference active HOR array, a broader alpha-satellite interval and the bases matching an HOR profile are different targets. Yq12 span and DYZ1-plus-DYZ2 unit bases also differ because the region contains smaller satellite classes and intervening sequence. Clear target definitions are needed both for calibration and for biological comparisons. Even quality-passing assemblies remain measurements with possible annotation and assembly errors; the reference checks reduce this uncertainty without eliminating it.

The smaller arrays require separate evaluation. Closely related repeat classes may share most candidate markers, leaving a small or biased set of diagnostic sequences. An aggregate result for Yq12 does not establish the accuracy of every component. Each output must be judged against its reference variability, held-out prediction error and relevant short or depleted states before it is used to infer population variation.

Sequencing protocol is a second limit. Counting all eligible read sequences removes dependence on their alignment positions for the matched files tested here, but it cannot restore sequence that a library or instrument underrepresents. Platform-specific scale factors correct average shifts, not arbitrary library-specific distortions. Restricted-region counting adds another condition: its calibration and capture efficiency depend on the alignment pipeline. These effects require validation at the level of the read group and analysis mode used in a cohort. Library G+C bias also remains within one project of PCR-free libraries. In \vsette{297} 1kGP parent--offspring pairs the parent--offspring difference followed the difference between the low-G+C coverage of the two libraries for DYZ18S (Spearman $\rho=\vsette{-0.509}$), DYZ1 (\vsette{$-0.393$}), DYZ17 (\vsette{$-0.382$}) and DYZ2 (\vsette{$+0.377$}), but not for DYZ3 (\vsette{$0.000$}) or DYZ19 (\vsette{$-0.126$}). Comparisons between males sequenced in different batches therefore carry extra error for these \vsette{four} classes. PCR-amplified libraries lie outside the calibration: in HGDP populations they shifted DYZ2 by \vsette{$-45.5\%$} and DYZ18S by \vsette{$+24.7\%$} against PCR-free libraries of the same populations.

Pedigrees illustrate how measurement bias can affect a biological interpretation. A small directional parent--offspring difference can arise when generations fall in different sequencing batches. Correcting such differences also requires care: using the same relatives to fit a batch adjustment and to test a transmission effect would couple calibration to the outcome. Independent calibration, explicit batch contrasts and family-aware uncertainty are needed before assigning a directional change to satellite biology. At the present percent-level resolution, pedigree concordance is most useful for detecting inconsistency; it does not measure the rate of small repeat-unit mutations.

The biological opportunity is to test the extent and exceptions of satellite compositional balance and to map supported array states across populations. Balanced DYZ1:DYZ2 abundance and their evolutionary coupling are already documented in the assembly resource \cite{Hallast2026}. A larger read-based catalog can examine reproducible departures, their distribution within lineages and their persistence across cohorts. The DYZ19 separation extends beyond the males with reference assemblies, survives close-kinship pruning and recurs in HGDP and SGDP; fine J sublineages and neighbouring groups remain to be tested.

These comparisons require a distinction between reference-panel observations and independent population evidence. The SNP tree supplies a phylogenetic display; remaining fine-scale topology and state-specific prediction uncertainty require verification before evolutionary inference. The tree has no time calibration. Measurement covariance, pedigree dependence and cohort effects must accompany models of satellite dynamics. A fitted tendency towards a common composition would not by itself demonstrate functional constraint or selection. Y dosage provides a parallel quality measurement for interpreting unusual array estimates.
